\documentclass[12pt,a4paper]{article}
% --- Pacotes Básicos ---
\usepackage[utf8]{inputenc}
\usepackage[T1]{fontenc}
%\usepackage[brazilian]{babel}
\usepackage{indentfirst}
\usepackage{geometry}
\geometry{left=3cm, top=3cm, right=2cm, bottom=2cm}

\usepackage{tabularx} % Para colunas com largura automática
\usepackage{array}    % Para centralizar texto verticalmente

% --- Informações do Documento ---
\title{Engenharia de Software II}
\author{Erickson Giesel Müller}
\date{\today}

\begin{document}

\maketitle

\section*{Definição do Tema}
\textbf{Tarefa:} Elicitação de Requisitos

\textbf{Descrição:} Definição dos requisitos funcionais e não-funcionais do Google Agenda.

%\section{Requisitos Não-Funcionais}
%\begin{table}[h]
%    \centering
%    \renewcommand{\arraystretch}{1.5} % Aumenta o espaçamento entre linhas
%    \begin{tabularx}{\textwidth}{|l|X|X|X|}
%        \hline
%        \textbf{ID} & \textbf{REQUISITO NÃO FUNCIONAL} & \textbf{DESCRIÇÃO DO REQUISITO NÃO FUNCIONAL}  \\ \hline
%        % Exemplos de linhas vazias para preenchimento:
%        RNF01        &                               &			\\ \hline
%        RNF02        &                               &			\\ \hline
%        RNF03        &                               &			\\ \hline
%    \end{tabularx}
%\end{table}

\section{Requisitos Funcionais}
\begin{table}[h]
    \centering
    \renewcommand{\arraystretch}{1.5} % Aumenta o espaçamento entre linhas
    \begin{tabularx}{\textwidth}{|l|X|X|X|}
        \hline
        \textbf{ID} & \textbf{REQUISITO FUNCIONAL} & \textbf{TIPO USUÁRIO DO SISTEMA} & \textbf{DESCRIÇÃO DO REQUISITO FUNCIONAL} \\ \hline
        % Exemplos de linhas vazias para preenchimento:
        RF01        &                               &                                  &                                           \\ \hline
        RF02        &                               &                                  &                                           \\ \hline
        RF03        &                               &                                  &                                           \\ \hline
    \end{tabularx}
\end{table}
\end{document}
